\documentclass[10pt,a4paper]{amsart}
\usepackage[margin=19mm]{geometry}
\usepackage{amsmath,amssymb,mathtools}
\usepackage[scr=rsfs,cal=euler]{mathalfa}
\usepackage{xfrac}
\usepackage[unicode,hidelinks,bookmarks=false]{hyperref}
\newcommand{\ie}{i.e.\ }
\newcommand{\cf}{cf.\ }
\newcommand{\resp}{respectively\ }
\newcommand{\Subsection}{\subsection}
\providecommand{\nobreakdash}{\nobreak\hbox{-}}
\setlength{\emergencystretch}{2em}
\multlinegap0pt
\usepackage{amssymb}
\usepackage{xcolor}
\usepackage{algorithm}\usepackage{algpseudocode}
\newtheorem{lemma}{Lemma}
\title{Fixed-Parameter Tractability of $t$-Uniform Hypergraphicality}
\author{Riley Brown, Istvan Miklos}
\date{Source: arXiv:2606.08523v1 (CC BY 4.0)}
\begin{document}
\maketitle

\noindent\emph{Source and licence.} Fixed-Parameter Tractability of $t$-Uniform Hypergraphicality. Version arXiv:2606.08523v1 (CC BY 4.0), distributed by arXiv under the Creative Commons Attribution 4.0 International licence, http://creativecommons.org/licenses/by/4.0/, as recorded in the arXiv licence metadata for this exact version and shown on the abstract page https://arxiv.org/abs/2606.08523v1. Full licence notice: \texttt{licenses/arxiv-2606.08523-CC-BY-4.0.txt}.\par\smallskip

\noindent\textit{Editorial scope.} Counted unit: the article's lemma lemma (source0.tex lines 291--293). The article's complete original proof of it is delivered with it (source0.tex lines 295--323, one \texttt{proof} environment of 29 lines). No mathematics is written, repaired or re-derived: the statement and the proof are the article's own lines, in the article's own notation, and no retained line is substituted or re-flowed.  No further source-local context is needed: every cross-reference of every retained line resolves inside the two delivered spans.  Nothing was omitted from any retained span, so the package carries no omission record.  No retained line contains a citation, so the wrapper carries no bibliography and no bibliography entry.  No retained line contains a figure environment, an image-inclusion command or drawing code, so no image is delivered and the package declares no asset dependency.  Editorial formatting only, in addition: an AMS \texttt{amsart} preamble that restates the article's own declarations and macros used by the retained lines, namely the environments \texttt{lemma} on separate counters, together with the standard packages those lines need.  The retained environments print in this document as Lemma 1 in order of appearance, whereas the article prints the same items with the numbering of its own full document.  The complete native source of the article as distributed in the arXiv e-print for this exact version is bundled unchanged as \texttt{sources/202335/source0.tex}.
\begin{lemma}[Existence of balancing hinge-flips]
Let $H$ be a $t$-uniform hypergraph and let $u,v \in V$ satisfy $d(u) > d(v)+1$. Then there exists a hinge-flip that decreases $d(u)$ by $1$ and increases $d(v)$ by $1$.
\end{lemma}

\begin{proof}
Let
\[
\mathcal{N}(w) = \{ S \subseteq V \setminus \{w\} : |S| = t-1,\ S \cup \{w\} \in E \},
\]
so that $|\mathcal{N}(w)| = d(w)$.

Partition $\mathcal{N}(u)$ into three parts:
\begin{align*}
\mathcal{A}_1 &= \{ S \in \mathcal{N}(u) : v \in S \}, \\
\mathcal{A}_2 &= \{ S \in \mathcal{N}(u) : v \notin S,\ S \in \mathcal{N}(v) \}, \\
\mathcal{A}_3 &= \{ S \in \mathcal{N}(u) : v \notin S,\ S \notin \mathcal{N}(v) \}.
\end{align*}

Thus $\mathcal{N}(u) = \mathcal{A}_1 \cup \mathcal{A}_2 \cup \mathcal{A}_3$ is a disjoint union.

We claim that
\[
|\mathcal{A}_1| + |\mathcal{A}_2| \le d(v).
\]
Indeed, each $S \in \mathcal{A}_1$ corresponds to an edge containing both $u$ and $v$, and each $S \in \mathcal{A}_2$ corresponds to an edge $S \cup \{v\}$ containing $v$ but not $u$. These edges are all distinct, hence their total number is at most $d(v)$.

Therefore,
\[
|\mathcal{A}_3| = d(u) - (|\mathcal{A}_1| + |\mathcal{A}_2|) \ge d(u) - d(v) > 0.
\]

Thus there exists $S \in \mathcal{A}_3$. Then $S \cup \{u\} \in E$, $v \notin S$, and $S \cup \{v\} \notin E$. Replacing $S \cup \{u\}$ by $S \cup \{v\}$ gives the desired hinge-flip.
\end{proof}

\end{document}
